\documentclass[11pt]{article}
\usepackage[margin=1in]{geometry}
\usepackage{amsmath}
\usepackage{graphicx}
\usepackage{hyperref}
\usepackage{booktabs}
\usepackage{float}
\usepackage{listings}
\usepackage{xcolor}
\usepackage{enumitem}

\definecolor{termbg}{gray}{0.95}
\definecolor{codebg}{rgb}{0.94,0.96,1.0}

\lstdefinestyle{termout}{
  basicstyle=\ttfamily\footnotesize,
  breaklines=true,
  breakatwhitespace=false,
  columns=fullflexible,
  frame=single,
  framerule=0.3pt,
  rulecolor=\color{black!40},
  backgroundcolor=\color{termbg},
  xleftmargin=4pt, xrightmargin=4pt,
  showstringspaces=false,
}

\lstdefinestyle{codeblk}{
  basicstyle=\ttfamily\footnotesize,
  breaklines=true,
  breakatwhitespace=false,
  columns=fullflexible,
  frame=single,
  framerule=0.3pt,
  rulecolor=\color{blue!40},
  backgroundcolor=\color{codebg},
  xleftmargin=4pt, xrightmargin=4pt,
  showstringspaces=false,
}

\title{Porting the UberClock LiteX SoC Build to openXC7:\\
Debugging Log, Fixes, and Status}
\author{}
\date{October 11, 2026}

\begin{document}
\maketitle

\begin{abstract}
\noindent
This note records an attempt to build the UberClock LiteX SoC
(\lstinline{2.soc-litex/}) for the Alinx AX7203 (\lstinline{xc7a200tfbg484-2})
using the fully open-source \textbf{openXC7} flow
(Yosys~$\to$~nextpnr-xilinx~$\to$~prjxray packer) instead of Vivado, as
tracked by item 7 of the project's own roadmap (``Port from Vivado to
openXC7''). It documents, in order, every error hit, its root cause, and
its fix (or current workaround). A minimal configuration --
RISC-V CPU, UART, integrated RAM, no DSP core, no DDR3, no Ethernet --
now builds and routes to a working \lstinline{.bit} file end to end. The
full design with the UberClock DSP core enabled
(\lstinline{--with-uberclock}) does \emph{not} yet build cleanly: it
exposes a reproducible nextpnr-xilinx placement-legality bug, independent
of the earlier capacity/performance issues already fixed. Section
\ref{sec:status} gives the current status and recommended path forward.
\end{abstract}

\tableofcontents

\section{Goal and starting point}

The project's existing Vivado-based LiteX build flow
(\lstinline{2.soc-litex/3.build/Makefile}, target script
\lstinline{8.python/src/targets/alinx_ax7203_uberclock.py}, SoC assembly in
\lstinline{8.python/src/uberclock_soc/soc.py}) already supports multiple
toolchains through LiteX's \lstinline{XilinxPlatform} abstraction. For
7-series devices LiteX lists four supported toolchains:
\lstinline{vivado}, \lstinline{f4pga}, \lstinline{yosys+nextpnr}, and
\lstinline{openxc7} (the latter two both route through
\lstinline{XilinxYosysNextpnrToolchain}, differing only in where chip
databases and the prjxray device database are expected to live). No RTL
changes were assumed to be necessary going in; the question was whether
the existing design would synthesize, place, and route on the open flow
at all, and if not, why.

\subsection{Toolchain environment discovered on this machine}

openXC7 was already installed, but not where LiteX's toolchain wrapper
defaults expect it (a snap path). The actual layout:

\begin{lstlisting}[style=termout]
/opt/openxc7/bin/            nextpnr-xilinx, yosys, bbasm, xc7frames2bit,
                              xc7patch, bitread, bittool, segmatch,
                              frame_address_decoder, gen_part_base_yaml
/opt/openxc7/share/prjxray-db/        device database (artix7/xc7a200t*)
/opt/openxc7/share/nextpnr-xilinx/python/   bbaexport.py (chipdb generator)
/opt/openxc7/setup_env.sh             ready-made env-var setup script
~/openxc7-chipdb/xc7a200tfbg484.bin   pre-generated chip database (317 MB)
\end{lstlisting}

The pre-generated chip database's filename was confirmed to match exactly
what LiteX's toolchain computes for this part. The toolchain derives a
\lstinline{dbpart} string from the device name via the regex
\lstinline{xc7([aksz])([0-9]+)(.*)-([0-9])} applied to
\lstinline{xc7a200tfbg484-2}, giving \lstinline{xc7a200tfbg484} --
an exact match to the cached file, so chip database generation (otherwise
a slow, memory-hungry step run through \lstinline{bbaexport.py} +
\lstinline{bbasm}) was skipped on every build in this session.

\section{Stage 1: CPU + UART + integrated RAM only}
\label{sec:stage1}

The first goal was the smallest possible configuration: RISC-V CPU, UART,
and the default integrated (BRAM-based) RAM, with no UberClock DSP core,
no UberDDR3, no Ethernet, no HDMI.

\subsection{Issue 1 -- Makefile recipe pasted into the shell instead of
invoked through \texttt{make}}

The first attempt failed before even reaching the toolchain:

\begin{lstlisting}[style=termout]
make build-board:
    mkdir -p $(CHIPDB)
    ...
make: *** No rule to make target 'build-board:'.  Stop.
zsh: command not found: CHIPDB
mkdir: missing operand
\end{lstlisting}

\textbf{Cause.} The Makefile recipe text had only been \emph{described}
in a prior chat turn, not actually written to the Makefile. It was pasted
directly into the interactive zsh prompt. zsh then (a) tried to run
\lstinline{make build-board:} literally, with a trailing colon, which is
not a valid target name, and (b) treated \lstinline{$(CHIPDB)} as shell
command substitution -- i.e. ``run a command literally named
\lstinline{CHIPDB} and substitute its output'' -- which obviously doesn't
exist as a command.

\textbf{Fix.} The toolchain-aware recipe was actually added to
\lstinline{2.soc-litex/3.build/Makefile} (see Appendix~\ref{app:makefile})
and invoked the correct way:
\begin{lstlisting}[style=termout]
make build-board TOOLCHAIN=openxc7 OPTIONS='...'
\end{lstlisting}

\subsection{Issue 2 -- \texttt{PRJXRAY\_DB\_DIR} unset in an interactive
shell}

A manual sanity check also failed:
\begin{lstlisting}[style=termout]
$ ls $PRJXRAY_DB_DIR/artix7/ | grep a200t
ls: cannot access '/artix7/': No such file or directory
\end{lstlisting}

\textbf{Cause.} \lstinline{PRJXRAY_DB_DIR} (and \lstinline{CHIPDB},
\lstinline{NEXTPNR_XILINX_PYTHON_DIR}) only existed inside the Makefile
recipe, never exported into the interactive shell.

\textbf{Fix.} Two options, both now valid: run
\lstinline{source /opt/openxc7/setup_env.sh} (a ready-made script already
present on disk), or simply go through \lstinline{make build-board
TOOLCHAIN=openxc7}, which now exports these variables itself (see
Appendix~\ref{app:makefile}).

\subsection{Issue 3 -- SA placer leaves a fracturable LUT pair unbound,
router then fails}

With the Makefile fixed and \lstinline{TOOLCHAIN=openxc7} building
correctly, the first real toolchain failure appeared during
place-and-route of the CPU-only design:

\begin{lstlisting}[style=termout]
Warning: fixupPlacement: slot with 7 unique inputs (>6): lut5=...
  lut6=...; skip
Warning: fixupPlacement: slot with 8 unique inputs (>6): lut5=...
  lut6=...; skip
  ... (150 such warnings total)
Info:     routing clock 's7mmcm0_clkout_buf'
Info:             failed to find a route using dedicated resources.
Info: Running router2...
ERROR: No wire found for port O6 on source cell
  $abc$27999$auto$blifparse.cc:536:parse_blif$28010.
150 warnings, 1 error
\end{lstlisting}

\textbf{Cause.} LiteX's \lstinline{XilinxYosysNextpnrToolchain} hardcodes
\lstinline{--placer sa} with no CLI override
(\lstinline{litex/build/xilinx/yosys_nextpnr.py}, constructor). On this
LUT-dense, fracturable-LUT-packed design (VexRiscv + its decode/ALU
logic), nextpnr-xilinx's SA placer runs a post-placement
\lstinline{fixupPlacement} pass that tries to legalize LUT5/LUT6 pairs
sharing one slice's six input pins. When it can't legalize a given pair
it just \lstinline{skip}s -- but this can leave one of the two cells
unbound to any BEL. The router then has no wire assignment for that
cell's output and fails.

\textbf{Fix.} Patched \lstinline{uberclock_soc/soc.py} to override the
placer right after platform construction:
\begin{lstlisting}[style=codeblk]
platform = alinx_ax7203.Platform(toolchain=toolchain)

if toolchain in ("openxc7", "yosys+nextpnr"):
    platform.toolchain._pnr_opts = platform.toolchain._pnr_opts.replace(
        "--placer sa", "--placer heap"
    )
\end{lstlisting}
Confirmed \lstinline{heap} is a supported placer
(\lstinline{nextpnr-xilinx --help} lists \lstinline{sa, heap}; the binary's
own undocumented default is actually \lstinline{heap}, but LiteX's wrapper
always passes \lstinline{sa} explicitly, overriding that default). The
\lstinline{heap} placer does not hit this particular legalization bug.

\subsection{Issue 4 -- BRAM CASCADE pair placed in non-adjacent columns}

With the placer fixed, a \emph{different} routing failure appeared:
\begin{lstlisting}[style=termout]
ERROR: Failed to route arc 0 of net 'main_ram.0.8.genblk1.genblk1.CAS_B',
  from SITEWIRE/RAMB36_X1Y35/CASCADEOUTB to SITEWIRE/RAMB36_X0Y45/CASCADEINB.
\end{lstlisting}

\textbf{Cause.} The design's default integrated RAM is 256~KiB
(\lstinline{INTEGRATED_MAIN_RAM_BYTES = 256 * 1024} in
\lstinline{BaseSoC}). At that depth, Yosys's BRAM inference uses a
depth-cascaded RAMB36E1 pair (dedicated \lstinline{CASCADEOUT}/
\lstinline{CASCADEIN} wiring, same-column-only resources on 7-series).
nextpnr-xilinx placed the two halves in \emph{different, non-adjacent
columns} (\lstinline{RAMB36_X1Y35} and \lstinline{RAMB36_X0Y45}) -- a
placement that is physically unroutable, because it is not enforcing the
cascade pair's required column adjacency as a relative placement
constraint.

\textbf{Fix (workaround, not a real fix).} Reduced the integrated RAM to
64~KiB for openXC7 builds via the already-exposed LiteX CLI flag:
\begin{lstlisting}[style=termout]
--integrated-main-ram-size=0x10000
\end{lstlisting}
small enough that Yosys no longer needs BRAM cascading for this memory.
\textbf{The 256~KiB default remains broken under openXC7.} A real fix
needs either an upstream nextpnr-xilinx fix to enforce CASCADE-pair
adjacency during placement, or a LiteX-side change to compose the RAM
differently (e.g.\ parallel narrower blocks addressed by higher bits)
for this toolchain.

\subsection{Issue 5 -- \texttt{fasm2frames} not found}

The final blocker for Stage~1:
\begin{lstlisting}[style=termout]
build_alinx_ax7203.sh: line 5: fasm2frames: command not found
\end{lstlisting}

\textbf{Cause.} The generated build script's packer stage is:
\begin{lstlisting}[style=termout]
fasm2frames --part xc7a200tfbg484-2 --db-root .../prjxray-db/artix7 \
  alinx_ax7203.fasm > alinx_ax7203.frames
xc7frames2bit --part_file .../part.yaml --part_name xc7a200tfbg484-2 \
  --frm_file alinx_ax7203.frames --output_file alinx_ax7203.bit
\end{lstlisting}
\lstinline{fasm2frames} is \emph{not} the \lstinline{fasm} PyPI package
(\lstinline{pip install fasm} only provides the \lstinline{fasm} text-format
CLI, a different tool, and was tried first on this mistaken assumption --
it installed fine but did not provide \lstinline{fasm2frames}). The real
\lstinline{fasm2frames} is a Python script,
\lstinline{utils/fasm2frames.py}, that lives in the
\href{https://github.com/f4pga/prjxray}{\texttt{f4pga/prjxray}}
repository itself, exposed as a console script via that repo's
\lstinline{setup.py}:
\begin{lstlisting}[style=codeblk]
entry_points={
    'console_scripts': ['fasm2frames=utils.fasm2frames:main'],
}
\end{lstlisting}
It is not published to PyPI under this name, and this machine's
\lstinline{/opt/openxc7} binary install does not ship it either.

\textbf{Fix.} Sparse-cloned just the needed subdirectories, fixed an
upstream packaging omission (\lstinline{utils} was never listed in
\lstinline{packages=[...]}, even though the entry point depends on it),
and installed it editable into the LiteX venv:
\begin{lstlisting}[style=termout]
git clone --depth 1 --filter=blob:none --sparse \
  https://github.com/f4pga/prjxray.git prjxray_src
cd prjxray_src && git sparse-checkout set utils prjxray
sed -i "s/packages=\['prjxray'\]/packages=['prjxray', 'utils']/" setup.py
pip install -e .
\end{lstlisting}
This also pulled in the real \lstinline{fasm}, \lstinline{intervaltree},
\lstinline{pyjson5}, \lstinline{pyyaml}, and \lstinline{simplejson}
dependencies declared in \lstinline{prjxray}'s \lstinline{setup.py}.
For reproducibility on a fresh machine, the same steps were added to
\lstinline{2.soc-litex/setup_dev.sh} behind a new \lstinline{--with-openxc7}
flag (Appendix~\ref{app:setupdev}).

\subsection{Stage 1 result: success}

\begin{lstlisting}[style=termout]
$ make build-board TOOLCHAIN=openxc7 \
    OPTIONS='--integrated-main-ram-size=0x10000 --nextpnr-seed=42'
...
exit=0

$ file alinx_ax7203.bit
alinx_ax7203.bit: Xilinx BIT data - from alinx_ax7203.frames;
  Generator=xc7frames2bit - for xc7a200tfbg484-2
  - built 2026/10/11(18:21:09) - data length 0x947a5c
\end{lstlisting}

A complete, working bitstream for RISC-V CPU~+~UART~+~64~KiB integrated
RAM, built entirely with Yosys~+~nextpnr-xilinx~+~prjxray, no Vivado
involved.

\section{Stage 2: enabling the UberClock DSP core}
\label{sec:stage2}

The next step added \lstinline{--with-uberclock}: the 11-channel DSP
pipeline (6~rx channels including a reference channel, 5~tx channels),
each instantiating a CORDIC stage chain plus CIC and half-band (HB)
decimation/interpolation filters (\lstinline{1.hw/uberclock/},
\lstinline{1.hw/cordic/}, \lstinline{1.hw/filters/}). DDR3 and Ethernet
remained disabled throughout Stage~2 -- those are expected to need Vivado
regardless (dropped \lstinline{INTERNAL_VREF} constraint, IDELAY/SERDES-
heavy PHYs, no real static timing analysis under this flow -- see
Section~\ref{sec:status}).

\subsection{Issue 6 -- \texttt{-flatten} before optimization turns a
multi-channel design into a multi-hour synthesis run}

The very first \lstinline{--with-uberclock} attempt never finished.
After nearly 10 minutes it was still deep inside Yosys's optimization
fixpoint loop, past its 26th round, with no error:
\begin{lstlisting}[style=termout]
26.14.31. Executing OPT_MUXTREE pass (detect dead branches in mux trees).
  ... <suppressed ~5669 debug messages>
26.14.32. Executing OPT_REDUCE pass (consolidate $*mux and $reduce_* inputs).
make: *** [Makefile:63: build-board] Terminated
\end{lstlisting}

\textbf{Cause.} LiteX's default synth options are
\lstinline{"-flatten -nosrl -arch xc7 "}. Yosys's own
\lstinline{synth_xilinx} documentation (\lstinline{yosys -p "help
synth_xilinx"}) confirms \lstinline{-flatten} is applied in the
\textbf{\lstinline{prepare:}} stage -- i.e.\ \emph{before}
\lstinline{map_dsp:}, \lstinline{coarse:}, \lstinline{map_memory:},
\lstinline{fine:}, and every \lstinline{opt_*} pass. Because the design
instantiates roughly 11 structurally identical rx/tx channel module
instances, flattening before optimization means Yosys re-runs its entire
(expensive, super-linear in places) optimization pipeline on the full
$\sim$11$\times$-larger flat netlist, instead of optimizing each unique
module once and simply copying the result 11 times at the end.

\textbf{Fix.} \lstinline{YosysNextPNRToolchain} exposes a documented,
first-class extension point for exactly this,
\lstinline{self._yosys_template} (default \lstinline{[]}, meaning ``use
the built-in default template''). Patched \lstinline{soc.py} to strip
\lstinline{-flatten} from \lstinline{_synth_opts} and supply a custom
template that keeps hierarchy through all of \lstinline{synth_xilinx}'s
internal optimization, flattening only immediately before the
\lstinline{write_json} step that nextpnr-xilinx's flat-netlist JSON input
requires:
\begin{lstlisting}[style=codeblk]
platform.toolchain._synth_opts = platform.toolchain._synth_opts.replace(
    "-flatten ", ""
)
platform.toolchain._yosys_template = [
    "verilog_defaults -push",
    "verilog_defaults -add -defer",
    "{read_files}",
    "verilog_defaults -pop",
    "attrmap -tocase keep -imap keep=\"true\" keep=1"
      " -imap keep=\"false\" keep=0 -remove keep=0",
    "{yosys_cmds}",
    "synth_{target} {synth_opts} -top {build_name}",
    "flatten",
    "clean",
    "write_{write_fmt} {write_opts} {output_name}.{synth_fmt}",
]
\end{lstlisting}
With this change, synthesis for the full 11-channel design actually
completes (where it previously did not, even given 30 minutes).

\subsection{Issue 7 -- without abc9, the mapped netlist is 117\% over LUT
capacity}

With the flatten fix alone (placer \lstinline{heap}, default Yosys
\lstinline{abc} mapper, 64~KiB RAM), synthesis completed and placement
began, then failed:
\begin{lstlisting}[style=termout]
Info: Creating initial analytic placement for 191206 cells, ...
ERROR: Failed to expand region (0, 0) |_> (264, 260) of 316493 SLICE_LUTXs
\end{lstlisting}
and nextpnr's own pre-placement utilization report showed:
\begin{lstlisting}[style=termout]
Info:     Created 314717 SLICE_LUTX cells from:
Info:          77172x FDRE
Info: 	          SLICE_LUTX: 316493/269200   117%
\end{lstlisting}

\textbf{Cause (confirmed, not assumed).} To rule out a real DSP-inference
regression from the Issue~6 fix, this was cross-checked against the
project's own \emph{known-working} Vivado synthesis report
(\lstinline{alinx_ax7203_utilization_synth.rpt}) for a config that
additionally includes the full UberDDR3 and Ethernet logic (i.e.\ a
\emph{larger} design than this openXC7 attempt):

\begin{table}[H]
\centering
\begin{tabular}{lrr}
\toprule
Resource & Vivado (DDR3+Eth+UberClock) & This run (UberClock only) \\
\midrule
DSP48E1   & 168            & 168 (exact match) \\
RAMB36E1  & 290            & 179 \\
RAMB18E1  & 22             & 19 \\
Slice LUTs / SLICE\_LUTX & 74448 / 134600 (55.3\%) & 316493 / 269200 (117\%) \\
\bottomrule
\end{tabular}
\caption{Vivado vs.\ this run's Yosys (\texttt{abc}, no \texttt{-abc9}) synthesis,
same UberClock RTL. (nextpnr-xilinx's \texttt{SLICE\_LUTX} capacity is reported
on a 2$\times$ finer basis than Vivado's whole-LUT6 \texttt{Slice LUTs} count, so
the two denominators differ by design, not by a counting bug.)}
\end{table}

The \lstinline{DSP48E1} count matching exactly (168 = 168) confirms DSP
inference itself was not the regression. The LUT blowup is specific to
Yosys's classic (non-\lstinline{abc9}) \lstinline{abc} LUT technology
mapper, which is well known to be markedly less efficient than Vivado's
proprietary mapper on this kind of RTL (deep CORDIC pipeline stages,
case-driven multiplexing).

\textbf{Fix.} Added \lstinline{--yosys-abc9} (already exposed on LiteX's
CLI via \lstinline{xilinx_yosys_nextpnr_args}, no code change required) to
switch to Yosys's newer \lstinline{abc9} LUT mapper. Result: total cell
count dropped from 191{,}206 to 131{,}503, and
\lstinline{SLICE_LUTX: 197705/269200 73%} -- comfortably within capacity
(DSP48E1 still exactly 168; RAMB36E1 179/365 = 49\%; RAMB18E1 19/730 =
2\%).

\subsection{Issue 8 -- HeAP placer cannot legalize an abc9-generated
CARRY4 chain (reproducible, not seed luck)}
\label{sec:issue8}

With \lstinline{abc9} added, placement progressed much further, then
failed differently:
\begin{lstlisting}[style=termout]
Info: Running main analytical placer.
ERROR: Unable to find legal placement for cell
  '$flatten\uberclock.\cordic_nco.\cordic_stages[16].stage_inst.
   $techmap$auto$opt_share.cc:199:merge_operators$2158173.
   $auto$alumacc.cc:512:replace_alu$2429962.genblk1.slice[0].genblk1.carry4',
  check constraints and utilisation.
\end{lstlisting}
At the time of this failure, overall utilization was comfortably within
budget (\lstinline{SLICE_LUTX} 73\%, \lstinline{SLICE_FFX} 39\%,
\lstinline{CARRY4} 48\%, \lstinline{RAMB36E1} 49\%, \lstinline{DSP48E1}
22\%) -- ruling out a capacity problem. This is a genuine
\textbf{placement-legality bug}: nextpnr-xilinx's HeAP placer's legalizer
cannot place this specific \lstinline{CARRY4} chain shape, which comes
from Yosys \lstinline{abc9}'s \lstinline{alumacc}/\lstinline{replace_alu}
ALU-extraction pass as applied to the CORDIC pipeline's adders.

To test whether this was placement-seed bad luck (as Issue~4's cascade
bug had been), the build was repeated with a different
\lstinline{--nextpnr-seed} (7 instead of 42), everything else identical.
Result: the \emph{same error signature}, at the \emph{same point} in
placement (immediately on entering ``Running main analytical placer''),
but on a \emph{different} \lstinline{CARRY4} instance
(\lstinline{cordic_cpu5.cordic_stages[14]} instead of
\lstinline{cordic_nco.cordic_stages[16]}):
\begin{lstlisting}[style=termout]
ERROR: Unable to find legal placement for cell
  '$flatten\uberclock.\cordic_cpu5.\cordic_stages[14].stage_inst...
   genblk1.carry4', check constraints and utilisation.
\end{lstlisting}
Two different seeds hitting the identical failure mode on two different
(but structurally identical) instances indicates the bug is
\textbf{systematic} -- HeAP cannot legalize \emph{any} CARRY4 of this
abc9-generated shape in this design -- not a probabilistic,
seed-dependent placement failure the way Issue~4 was. Further seed
rolling was judged unlikely to help and was not pursued further.

\subsection{Issue 9 -- disabling CARRY4 does not fix the bug, it just
relocates it, non-deterministically}

As a workaround for Issue~\ref{sec:issue8}, \lstinline{-nocarry} was
added to \lstinline{synth_xilinx}'s options (forcing the CORDIC adders
into plain LUT logic instead of dedicated \lstinline{CARRY4} chains),
keeping \lstinline{abc9} and the \lstinline{heap} placer:
\begin{lstlisting}[style=codeblk]
if os.environ.get("OPENXC7_NOCARRY"):
    platform.toolchain._synth_opts += "-nocarry "
\end{lstlisting}

\textbf{First attempt (30-minute cap).} The \lstinline{CARRY4} crash
disappeared (\lstinline{CARRY4: 0/33650 0%}), LUT utilization remained
well within budget (\lstinline{SLICE_LUTX: 152384/269200 56%}), total
placer cell count grew to 204{,}521 (from 131{,}503 with
\lstinline{CARRY4} enabled), and the HeAP placer printed \emph{no}
iteration output at all before the 30-minute cap was hit:
\begin{lstlisting}[style=termout]
Info: Running main analytical placer.
make: *** [Makefile:63: build-board] Terminated
\end{lstlisting}
Taken alone, this looked like non-convergence rather than a crash.

\textbf{Second attempt, same command, no premature cap.} To find out
whether that was genuine non-convergence or just a slow instance of the
same failure, the identical command was re-run with no artificial time
limit (a 4-hour safety net instead of a 30-minute cutoff). This time it
did \emph{not} hang: it reached the identical point --
``\texttt{Running main analytical placer.}'' -- and failed on the very
next line, almost instantly, with the same class of error as Issue~8 but
on a \emph{different} cell, this time one produced by Yosys
\lstinline{abc9}'s own internal reintegration pass rather than the CORDIC
adder's \lstinline{CARRY4}:
\begin{lstlisting}[style=termout]
Info: Running main analytical placer.
ERROR: Unable to find legal placement for cell
  '$flatten\uberclock.$auto$abc9_ops.cc:1552:reintegrate$3585864',
  check constraints and utilisation.
\end{lstlisting}
As before, utilization at the point of failure was comfortably within
budget (\lstinline{SLICE_LUTX: 152441/269200 56%}, \lstinline{CARRY4:
0/33650 0%}) -- not a capacity problem -- and, as in both Issue~8 cases,
\emph{zero} ``at iteration \#1...'' lines were ever printed; the failure
happens before annealing proper even starts.

\textbf{Revised conclusion.} Two runs of the \emph{exact same command}
(same seed, same options) produced different symptoms -- one an
apparent 30-minute hang with no output, the other an almost-instant
crash -- at the identical point in placement. \lstinline{-nocarry} does
not fix Issue~\ref{sec:issue8}'s bug; it relocates the same underlying
nextpnr-xilinx placer weakness (failing to legalize a specific
abc9-produced cell during HeAP's very first legalization step, before
any annealing) from a \lstinline{CARRY4} chain onto a different cell
type, and whether that manifests as a quick, explicit error or an
unbounded-looking stall varies non-deterministically between otherwise
identical runs. This is a stronger finding than ``disabling CARRY4
regresses convergence'': the bug is not specific to \lstinline{CARRY4}
at all, and no synthesis-option change tried so far removes it, only
moves it.

\textbf{Conclusion at this point.} \lstinline{CARRY4} enabled
$\Rightarrow$ placer fails to legalize a \lstinline{CARRY4} cell
(Issue~8). \lstinline{CARRY4} disabled $\Rightarrow$ placer fails to
legalize a different abc9-internal cell instead, non-deterministically
presenting as either a crash or an apparent hang (Issue~9). Both paths
hit the same class of nextpnr-xilinx placer bug; neither is a real fix.

\subsection{Issue 10 -- SA placer with abc9: completes annealing, then
hits a third, independent nextpnr-xilinx legalization bug}
\label{sec:issue10}

Keeping \lstinline{abc9} and \lstinline{CARRY4} enabled (i.e.\ without
\lstinline{-nocarry}), the placer was switched back to \lstinline{sa} --
a combination not yet tried, since the original \lstinline{sa} failure
(Issue~3) was observed on a much smaller, non-\lstinline{abc9} CPU-only
netlist with a completely different cell mix:
\begin{lstlisting}[style=termout]
OPENXC7_PLACER=sa \
make build-board TOOLCHAIN=openxc7 \
  OPTIONS='--with-uberclock --integrated-main-ram-size=0x10000 \
           --nextpnr-seed=42 --yosys-abc9'
\end{lstlisting}
(enabled via a new, temporary \lstinline{OPENXC7_PLACER} environment-variable
override in \lstinline{soc.py}, see Appendix~\ref{app:socpy}).

\textbf{First attempt (30-minute cap).} No crash and no placement-legality
error were observed, but the run was only $\sim$19\% through nextpnr's
very first, non-optimizing placement pass when the cap was hit (see the
original draft of this section for the exact trace) -- inconclusive on
whether \lstinline{sa} would eventually also hit Issue~8's bug class.

\textbf{Second attempt, no premature cap (6-hour safety net).} This time
the run went much further: it completed \lstinline{sa}'s entire
annealing schedule -- 272 iterations, converging on a stable timing cost,
over 5527 seconds ($\approx$92 minutes) of actual SA placement time --
and only then failed, during post-placement legalization:
\begin{lstlisting}[style=termout]
Info:   at iteration #272: temp = 0.000000, timing cost = 11993, wirelen = 71860986
Info: SA placement time 5527.12s
ERROR: post-placement validity check failed for Bel 'SLICE_X0Y249/A5FF' (no cell)
ERROR: Placing design failed.
\end{lstlisting}
Utilization at the point of failure was identical to Issue~8's
\lstinline{heap} run on the same netlist (\lstinline{SLICE_LUTX}
197707/269200 = 73\%, \lstinline{CARRY4} 16209/33650 = 48\%) -- again
not a capacity problem. nextpnr's own timing analysis also flagged 64
combinational loops in the netlist throughout the run, explicitly
attributed to \lstinline{abc9} reintegration and treated as non-fatal by
default; this did not cause the failure but is noted here as a property
of the \lstinline{abc9}-mapped netlist worth tracking if timing closure
is attempted later.

\textbf{This is a third, independently reproducible bug signature.} The
exact message \lstinline{"post-placement validity check failed for Bel
'<slice>/A5FF' (no cell)"} was also seen, on a \emph{different, earlier}
netlist, during Stage~1 seed-rolling: the very first attempt to use the
\lstinline{heap} placer on the (pre-\lstinline{abc9}) 256~KiB-RAM
CPU-only build hit \lstinline{"SLICE_X10Y181/A5FF"} in exactly the same
way, when \lstinline{heap}'s own brief internal SA-refinement sub-step
ran after its analytic placement. That occurrence was resolved along the
way to Stage~1's working configuration (Section~\ref{sec:stage1}) and
was not written up as its own numbered issue there; it is referenced here
because the same bug reappearing, independently, on a completely
different (larger, \lstinline{abc9}-mapped) netlist confirms it is a
general nextpnr-xilinx SA-legalization weakness, not an artifact specific
to one design size or one synthesis configuration.

\textbf{Conclusion.} \lstinline{sa} does not avoid Issue~8's problem --
it runs substantially longer (a full, successful anneal) and then hits a
\emph{different} systematic nextpnr-xilinx legalization bug instead. Across
Issues~8--10, every placer/synthesis-option combination tried
(\lstinline{heap}; \lstinline{heap}+\lstinline{-nocarry}; \lstinline{sa})
on the full UberClock netlist has hit an independent placement-legality
bug, and one of those bug signatures (\lstinline{A5FF (no cell)}) recurs
across two unrelated netlist configurations. This is strong evidence of
a genuine, multi-faceted toolchain maturity gap rather than a single
fixable misconfiguration on this project's side.

\section{Summary of attempts}

\begin{table}[H]
\centering
\small
\begin{tabular}{p{0.20\textwidth}p{0.30\textwidth}p{0.15\textwidth}p{0.25\textwidth}}
\toprule
Configuration & Key toolchain options & Result & Failure / outcome \\
\midrule
CPU+UART+RAM & \lstinline{sa} placer (LiteX default) & Fail & LUT5/6 fixupPlacement $\to$ unroutable (Issue~3) \\
CPU+UART+RAM & \lstinline{heap} placer & Fail & BRAM CASCADE non-adjacent (Issue~4) \\
CPU+UART+64KiB RAM & \lstinline{heap} placer & \textbf{Success} & Working \lstinline{.bit} \\
+UberClock & \lstinline{heap}, \lstinline{-flatten} (default) & Fail (no error) & Non-converging synth, 10+ min, round 26 of opt loop (Issue~6) \\
+UberClock & \lstinline{heap}, hierarchy-preserving flatten & Fail & 117\% LUT overcapacity (Issue~7) \\
+UberClock & \lstinline{heap}, +\lstinline{abc9} & Fail & CARRY4 placement-legality bug, seed 42 (Issue~8) \\
+UberClock & \lstinline{heap}, +\lstinline{abc9}, seed 7 & Fail & Same bug class, different CARRY4 instance (Issue~8) \\
+UberClock & \lstinline{heap}, +\lstinline{abc9}, +\lstinline{-nocarry} & Fail (2$\times$) & Same bug class on a different cell; hang vs.\ instant crash, non-deterministic (Issue~9) \\
+UberClock & \lstinline{sa}, +\lstinline{abc9} & Fail & Full 92-min anneal converges, then \texttt{A5FF (no cell)} legalization bug -- 3rd distinct bug signature (Issue~10) \\
\bottomrule
\end{tabular}
\caption{All configurations attempted, in chronological order.}
\end{table}

\section{Current status and recommended path forward}
\label{sec:status}

\subsection{What works today}

A complete, Vivado-free build for RISC-V CPU~+~UART~+~64~KiB integrated
RAM on the Alinx AX7203, via:
\begin{lstlisting}[style=termout]
make build-board TOOLCHAIN=openxc7 \
  OPTIONS='--integrated-main-ram-size=0x10000 --nextpnr-seed=42'
\end{lstlisting}
This exercises the full open flow -- Yosys, nextpnr-xilinx, prjxray's
\lstinline{fasm2frames} and \lstinline{xc7frames2bit} -- end to end, and
is loadable via the existing \lstinline{OpenFPGALoader}-based programmer
(\lstinline{alinx_ax7203.py}'s \lstinline{create_programmer}, unchanged by
this work).

\subsection{What does not work today}

\lstinline{--with-uberclock} (the DSP core) does not yet build cleanly
under openXC7. The blocking issue is \emph{not} capacity -- it is
nextpnr-xilinx placement-legalization bugs, and testing has now turned up
\textbf{three independent ones} on this one design:
\begin{itemize}[nosep]
  \item \lstinline{heap} fails to legalize a specific \lstinline{abc9}-
  produced cell during its very first legalization step, before any
  annealing runs -- first seen on a \lstinline{CARRY4} chain from the
  CORDIC adders (Section~\ref{sec:issue8}), reproducible across random
  seeds.
  \item Disabling \lstinline{CARRY4} inference (Issue~9) does not fix
  that bug, it relocates it onto a different \lstinline{abc9}-internal
  cell, presenting non-deterministically as either an instant error or
  an apparent multi-hour stall across otherwise-identical runs.
  \item Switching to the \lstinline{sa} placer (Issue~10) avoids both of
  the above -- it completes a full, successful 92-minute anneal -- but
  then hits a \emph{third}, independent bug during post-placement
  legalization (\lstinline{"post-placement validity check failed for Bel
  '<slice>/A5FF' (no cell)"}), the same exact signature already seen once
  before, on an unrelated netlist, during Stage~1 seed-rolling.
\end{itemize}
No placer/synthesis-option combination tried gets past placement for the
full design, and one bug signature (\lstinline{A5FF (no cell)}) has now
been independently reproduced on two different netlists, which makes it
the single strongest upstream bug-report candidate of the three.

\lstinline{--with-uberddr3}, \lstinline{--with-ethernet}, and
\lstinline{--with-hdmi} were not attempted under openXC7 in this session
and are expected to need Vivado regardless, independent of the issues
above:
\begin{itemize}[nosep]
  \item UberDDR3's \lstinline{INTERNAL_VREF} and
  \lstinline{BITSTREAM.STARTUP.MATCH_CYCLE} Vivado properties are
  silently dropped by \lstinline{XilinxPlatform.add_platform_command}
  for this toolchain (with a printed warning) -- without
  \lstinline{INTERNAL_VREF}, the DDR3 I/O bank has no reference level.
  \item The DDR3 PHY instantiates 27$\times$\lstinline{OSERDESE2},
  12$\times$\lstinline{ISERDESE2}, 9$\times$\lstinline{IDELAYCTRL},
  13$\times$\lstinline{IDELAYE2}, in \lstinline{VAR_LOAD} calibration
  mode -- the least-exercised part of nextpnr-xilinx's primitive support.
  \item \lstinline{add_false_path_constraint} is a documented no-op in
  LiteX's \lstinline{XilinxYosysNextpnrToolchain}; nextpnr-xilinx derives
  only a single global \lstinline{--freq} from the highest
  \lstinline{create_clock}, with no real propagation through the MMCM --
  so the DDR3 PHY's 400~MHz domain is effectively unconstrained and
  unverified under this flow regardless of the CARRY4 issue.
\end{itemize}

\subsection{Recommended next steps, in order of effort}

\begin{enumerate}
  \item \textbf{Treat all three legalization bugs (Issues~8--10) as
  genuine upstream nextpnr-xilinx limitations}, not something fixable
  from this project's RTL or LiteX integration. Every placer/option
  combination tried hit one: \lstinline{heap} hits the \lstinline{abc9}-
  cell bug directly; \lstinline{-nocarry} relocates it; \lstinline{sa}
  avoids both but hits the independent \lstinline{A5FF (no cell)} bug
  after a full, successful anneal. Prioritize filing the
  \lstinline{A5FF (no cell)} bug first -- it is the one already
  reproduced independently on two unrelated netlists (Stage~1's
  pre-\lstinline{abc9} CPU-only build and Stage~2's full
  \lstinline{abc9}-mapped design), the strongest signal of a general,
  fixable toolchain defect rather than a one-off. File minimal
  reproducers (a single CORDIC stage module is likely enough for the
  \lstinline{abc9}-cell bug) against
  \href{https://github.com/openXC7/toolchain-installer}{openXC7} or
  \href{https://github.com/gatecat/nextpnr-xilinx}{nextpnr-xilinx}
  upstream, and check whether a newer toolchain build already fixes
  either.
  \item \textbf{In the meantime}, keep Vivado as the synthesis path for
  any configuration that includes \lstinline{--with-uberclock},
  \lstinline{--with-uberddr3}, \lstinline{--with-ethernet}, or
  \lstinline{--with-hdmi}. Use openXC7 for the CPU-only bring-up
  configuration documented above, and for any future peripheral-free RTL
  bring-up work where a fully open flow is valuable.
  \item \textbf{Longer term, if pursuing openXC7 further}: consider
  whether avoiding \lstinline{abc9}'s \lstinline{alumacc}/
  \lstinline{replace_alu} ALU-extraction specifically (rather than
  disabling \lstinline{CARRY4} globally with \lstinline{-nocarry}, which
  regressed placer convergence) is possible -- e.g.\ by restructuring the
  CORDIC stage's adders in RTL to a shape Yosys maps more conventionally,
  or by scripting a custom Yosys pass ordering that skips just that
  substitution. This is invasive and was not attempted in this session.
\end{enumerate}

\section*{Appendix}
\appendix

\section{Makefile changes}
\label{app:makefile}

Added to \lstinline{2.soc-litex/3.build/Makefile}:

\begin{lstlisting}[style=codeblk]
# Toolchain: vivado (default) or openxc7
TOOLCHAIN     ?= vivado

# openXC7 install (actual layout on this machine: /opt/openxc7, not snap)
OPENXC7_DIR               ?= /opt/openxc7
CHIPDB                    ?= $(HOME)/openxc7-chipdb
PRJXRAY_DB_DIR            ?= $(OPENXC7_DIR)/share/prjxray-db
NEXTPNR_XILINX_PYTHON_DIR ?= $(OPENXC7_DIR)/share/nextpnr-xilinx/python
\end{lstlisting}

and the \lstinline{build-board} recipe:
\begin{lstlisting}[style=codeblk]
build-board:
	@echo "[build] target=$(TARGET_PY)"
	@echo "[build] toolchain=$(TOOLCHAIN) sys_clk=$(FREQ) options=$(OPTIONS)"
	source $(LITEX_VENV)/bin/activate && \
	if [ "$(TOOLCHAIN)" = "openxc7" ]; then \
	  mkdir -p "$(CHIPDB)"; \
	  export PATH="$(OPENXC7_DIR)/bin:$$PATH"; \
	  export CHIPDB="$(CHIPDB)"; \
	  export PRJXRAY_DB_DIR="$(PRJXRAY_DB_DIR)"; \
	  export NEXTPNR_XILINX_PYTHON_DIR="$(NEXTPNR_XILINX_PYTHON_DIR)"; \
	fi; \
	PYTHONPATH=$(PYTHONPATH_REPO) \
	$(PYTHON) $(TARGET_PY) \
	  --build \
	  --toolchain=$(TOOLCHAIN) \
	  --sys-clk-freq=$(FREQ) \
	  $(OPTIONS)
\end{lstlisting}

\section{\texttt{soc.py} changes}
\label{app:socpy}

Added to \lstinline{8.python/src/uberclock_soc/soc.py}, immediately after
platform construction in \lstinline{BaseSoC.__init__}:

\begin{lstlisting}[style=codeblk]
platform = alinx_ax7203.Platform(toolchain=toolchain)

if toolchain in ("openxc7", "yosys+nextpnr"):
    # Issue 3 / Issue 10: placer choice, overridable while bisecting
    # nextpnr-xilinx placement-legality bugs (sa: fracturable-LUT-pair
    # fixupPlacement bug; heap: CARRY4 legalization bug on abc9's
    # alumacc/replace_alu shapes). TEMPORARY env-gated override.
    _placer = os.environ.get("OPENXC7_PLACER", "heap")
    platform.toolchain._pnr_opts = platform.toolchain._pnr_opts.replace(
        "--placer sa", f"--placer {_placer}"
    )

    # Issue 6: keep hierarchy through synth_xilinx (don't -flatten before
    # optimization), flatten only right before nextpnr's required-flat
    # JSON output.
    platform.toolchain._synth_opts = platform.toolchain._synth_opts.replace(
        "-flatten ", ""
    )
    platform.toolchain._yosys_template = [
        "verilog_defaults -push",
        "verilog_defaults -add -defer",
        "{read_files}",
        "verilog_defaults -pop",
        "attrmap -tocase keep -imap keep=\"true\" keep=1"
          " -imap keep=\"false\" keep=0 -remove keep=0",
        "{yosys_cmds}",
        "synth_{target} {synth_opts} -top {build_name}",
        "flatten",
        "clean",
        "write_{write_fmt} {write_opts} {output_name}.{synth_fmt}",
    ]

    # Issue 9: TEMPORARY env-gated experiment, disables CARRY4 inference
    # to sidestep Issue 8's placement-legality bug. Regressed placer
    # convergence instead (see Issue 9) -- not enabled by default.
    if os.environ.get("OPENXC7_NOCARRY"):
        platform.toolchain._synth_opts += "-nocarry "
\end{lstlisting}

\noindent\textbf{Note:} the \lstinline{OPENXC7_PLACER} and
\lstinline{OPENXC7_NOCARRY} environment-variable gates are deliberately
temporary, added to bisect the Stage~2 failures without committing to a
default before knowing which combination (if any) actually works. Once
Issue~10 resolves, these should be collapsed into either a fixed default
or a proper \lstinline{--yosys-nocarry}-style CLI flag, not left as
undocumented environment variables.

\section{\texttt{setup\_dev.sh} changes}
\label{app:setupdev}

Added a \lstinline{--with-openxc7} flag to
\lstinline{2.soc-litex/setup_dev.sh} that performs the same
\lstinline{fasm2frames} install as Issue~5's manual fix, so a fresh
machine doesn't have to rediscover it:

\begin{lstlisting}[style=codeblk]
if [[ "${WITH_OPENXC7}" -eq 1 ]]; then
  echo "--- Installing fasm2frames (prjxray) for openXC7 ---"
  PRJXRAY_SRC_DIR="${LITEX_DIR}/prjxray_fasm2frames"
  if ! python -c "import utils.fasm2frames" >/dev/null 2>&1; then
    rm -rf "${PRJXRAY_SRC_DIR}"
    git clone --depth 1 --filter=blob:none --sparse \
      https://github.com/f4pga/prjxray.git "${PRJXRAY_SRC_DIR}"
    (cd "${PRJXRAY_SRC_DIR}" && git sparse-checkout set utils prjxray)
    sed -i "s/packages=\['prjxray'\]/packages=['prjxray', 'utils']/" \
      "${PRJXRAY_SRC_DIR}/setup.py"
    pip install -e "${PRJXRAY_SRC_DIR}"
  fi
  command -v fasm2frames >/dev/null 2>&1 \
    && echo "fasm2frames OK: $(command -v fasm2frames)" \
    || echo "WARNING: fasm2frames still not on PATH after install."
fi
\end{lstlisting}

\end{document}
